\chapter{Allgemeines}
\label{chap:allgemeines}

\section{Aufbau der Arbeit}
Der grundsätzliche Aufbau entspricht dem dieses Dokuments, so dass Sie dieses auch für Ihre Arbeit verwenden können:
\begin{enumerate}
	\item Deckblatt
	\item Kopie der Originalaufgabenstellung
	\item Ehrenwörtliche Erklärung selbständiger Bearbeitung
	\item verschiedene Verzeichnisse (je nach Bedarf)
	\item Kurzzusammenfassung deutsch/englisch (falls vom Betreuer gefordert)
	\item Einleitung
	\item Hauptteil
	\begin{enumerate}
		\item Stand der Technik
		\item Konzept
		\item Implementierung
		\item Ergebnisse 
	\end{enumerate}
	\item Zusammenfassung und Ausblick
	\item Literaturverzeichnis
	\item Anhang (optional)
\end{enumerate}

Diese Gliederung ist sehr umfassend und gilt nur als Anhalt. Insbesondere die Anzahl der verwende-ten Verzeichnisse ist insbesondere von der Aufgabenstellung.

\section{Kurzzusammenfassung (Abstract)}
Von der Arbeit ist eine Kurzzusammenfassung (Abstract) anzufertigen. Diese umfasst auf maximal einer Seite eine generelle Zusammenfassung der Arbeit inklusive Aufgabenstellung sowie die wesent-lichen Ergebnisse. Nach Absprache mit dem Betreuer auch in englischer Form.

\section{Quellenangaben und Zitate}
Jede Verwendung fremden geistigen Eigentums ist durch eine genaue Quellenangabe kenntlich zu machen. Die Kenntlichmachung erfolgt, indem dem jeweiligen Satz oder Abschnitt ein Verweis auf die zugehörige Literaturstelle angefügt wird. Dies erfolgt in der Form „[Krü 98]“, wobei die ersten 3 Buch-staben des Nachnamens und das Erscheinungsjahr verwendet werden.
„Wörtliche Zitate sind kursiv in Anführungszeichen zu schreiben“. Sinngemäß übernommene Textpas-sagen bzw. Aussagen werden nicht in Anführungsstrichen gesetzt.

\section{Schreibstil}
In der Regel ist ein unpersönlicher Schreibstil zu verwenden (nicht: „Ich habe…“). Angeführte Argumente sind logisch zu begründen oder mit Quellenverweisen zu belegen und zu untermauern. Auf persönliche Werturteile ist weitgehend zu verzichten. Durchgeführte Arbeiten sind in der Vergangenheitsform zu beschreiben, allgemeingültige Erkenntnisse in der Gegenwartsform. Beispiel: „Es wurde eine Formel entwickelt, welche das Verhalten gut beschreibt.“

\section{Quellenangaben und Zitate}
Jede Verwendung fremden geistigen Eigentums ist durch eine genaue Quellenangabe kenntlich zu machen. Die Kenntlichmachung erfolgt, indem dem jeweiligen Satz oder Abschnitt ein Verweis auf die zugehörige Literaturquelle angefügt wird. Dies erfolgt in der Form "\cite{Hal_SicherheitsgerichteteEchtzeitsysteme}", wobei die ersten drei Buchstaben des Nachnamens und das Erscheinungsjahr verwendet werden. Die Seitenzahl muss nicht angegeben werden. Beispiele hierzu können Sie dem beispielhaften Literaturverzeichnis dieses Dokuments entnehmen. In dieser Vorlage ist der korrekte Zitationsstil bereits eingestellt, sodass Sie die Quellen lediglich in den korrekten BibTex-Dateien einsortieren müssen. Achten Sie außerdem darauf, aus Ihrem Literaturprogramm (bspw. Citaiv) den entsprechenden sog. Kurzbeleg zu exportieren. Für Citavi wird ein entsprechender Kurzbeleg angeboten. Sprechen Sie Ihren Betreuer auf das Einrichten des Kurzbelegs und den korrekten Export aus Ihrer Literatursoftware an. 
„Wörtliche Zitate sind kursiv in Anführungszeichen zu schreiben“. Sinngemäß übernommene Textpassagen bzw. Aussagen werden nicht in Anführungsstrichen gesetzt.
Sind Verweise auf die Literaturstelle sehr lang, so ist es möglich, diese in die Fußzeile zu verschieben.

